Koncentrace majetku na~vrcholu rozdělení je odhadována pomocí distribučních národních účtů v~kombinaci s~daňovými záznamy a~firemními registry, jak je shromažďuje databáze \acl{WID}.
Na~rozdíl od~průzkumů domácností zachycuje \acs{WID} také skrytý kapitál holdingů, zahraničních fondů a~dalších finančních nástrojů, čímž lépe podchycuje horní zlomek procenta vlastníků.
Mapové znázornění zachycuje podíl nejbohatšího~\SI{1}{\percent} dospělých na~celkovém čistém jmění domácností za~rok~2024.

\acloc{ČR} vykazuje podíl horního~\SI{1}{\percent} na~hodnotě~\SI{26{,}3}{\percent}, přičemž mediánová hodnota v~rámci~\acs{EU}-27 čítá~\SI{24{,}4}{\percent}.
\acs{ČR} se tak řadí na~17.~místo z~27 členských států --- koncentrace je vyšší, než u~\SI{60}{\percent} zemí~\acp{EU}.
Klíčovým kontextuálním paradoxem je asymetrie mezi příjmovou a~majetkovou nerovností: příjmová nerovnost v~\acloc{ČR} (Giniho koeficient příjmů~\SI{27}{\percent} dle Eurostatu) patří k~nejnižším v~\acloc{EU}, zatímco koncentrace majetku na~vrcholu je nadprůměrná.
Tato divergence odráží dědictví privatizace po~roce~1989, kdy se průmyslové a~obchodní holdingy soustředily v~rukou relativně malé skupiny vlastníků kuponové privatizace a~pozdějších konsolidací, přičemž mzdová distribuce zůstala po~desetiletí plochá díky kolektivním dohodám i~strukturálním rysům českého pracovního trhu.

V~rámci~\acp{EU} je variabilita značná: od~\SI{14{,}0}{\percent} v~Nizozemsku přes~\SI{17{,}2}{\percent} na~Slovensku a~\SI{18{,}0}{\percent} ve~Finsku až~po~\SI{35{,}0}{\percent} v~Estonsku a~\SI{33{,}4}{\percent} v~Maďarsku.
Nižší koncentraci vykazují zejména severské země a~Benelux, v~nichž silný \acl{KV} a~rozsáhlé veřejné penzijní fondy distribuují výnosy z~kapitálu šíře než systémy postavené na~individuálním spoření.
V~postkomunistickém prostoru je rozptyl rovněž výrazný: zatímco Slovensko a~Slovinsko (\SI{23{,}2}{\percent}) vykazují nižší koncentraci než \acs{ČR}, Polsko (\SI{29{,}9}{\percent}), Lotyšsko (\SI{28{,}6}{\percent}) a~pobaltské státy jsou výrazně vyšší, což odráží odlišné trajektorie privatizace a~míru oligarchizace v~jednotlivých tranzičních ekonomikách.
Srovnání \aclins{ČR} a~Slovenska je zvláště ilustrativní: nižší slovenský podíl odpovídá přímějšímu vstupu zahraničních investorů s~atomizovanou vlastnickou strukturou (automobilový průmysl), zatímco v~\acloc{ČR} se průmyslové holdingové skupiny konsolidovaly v~rukou domácích vlastníků.

Hlubší kontext pro interpretaci majetkové koncentrace nabízí teze o~takzvaném technofeudalismu.
Yanis Varoufakis~\cite{varoufakis_technofeudalism_2023} argumentuje, že digitální platformy (Amazon, Google, Meta a~další) přetvořily tržní kapitalismus na~systém extrakce renty analogický feudálnímu uspořádání: vlastníci cloudové infrastruktury --- "cloudalisté" --- inkasují rentu od~prodejců, řidičů a~tvůrců obsahu, kteří jsou na~platformách závislí jako "clouďoví nevolníci".
Tato transformace má přímý dopad na~kolektivní vyjednávání: gig-ekonomika a~platformová práce nepodléhají tradičním pracovněprávním vztahům, pracovníci jsou klasifikováni jako \acp{OSVČ} bez nároku na~odborové zastoupení, a~jejich vyjednávací pozice je systematicky oslabena algoritmickým managementem a~snadností nahradit je dalšími "nevolníky".
Digitální kapitál koncentrovaný v~hrstce globálních platforem vstupuje do~národních statistik bohatství jen nedokonale: \ac{WID} nepostihuje plně tržní kapitalizaci nesvěřeneckých holdingů a~offshore struktur, jimiž je digitální majetek obvykle organizován, a~skutečná koncentrace může být tedy v~\acloc{ČR} i~v~celoevropském měřítku vyšší, než zobrazená mapa naznačuje.
V~kontextu diskuse o~\aclploc{KS} a~mzdové politice v~\acloc{ČR} platí: pokud část pracovní síly přechází z~regulovaného zaměstnaneckého poměru do~platformového modelu, klesá celkové pokrytí \acs{KS} a~tím i~faktická vyjednávací síla zbývajících odborových organizací.

% Editable figure definition for eu_bohatstvi_top1_wid
% Edit freely --- Python will NOT overwrite this file.
% To regenerate defaults, delete this file and re-run the script.
\def\strEuBohatstviTopIWidTitle{Podíl top 1\,\% domácností na čistém jmění (WID, do 2024)}%
\def\strEuBohatstviTopIWidColorbarLabel{podíl top 1\,\% na čistém jmění [\%]}%
\def\strEuBohatstviTopIWidCaption{Podíl top 1\,\% domácností na čistém jmění, mapa \acs{EU}, do~2024. Zdroj dat: \acs{WID}~\cite{wid_world_shweal}.}%
%

% --- Auto-added y-nudge knobs (override with \renewcommand)
\providecommand\NudgeEuBohatstviTopIWidCz{0pt}%
\providecommand\NudgeEuBohatstviTopIWidDk{0pt}%
\providecommand\NudgeEuBohatstviTopIWidAt{6pt}%
\providecommand\NudgeEuBohatstviTopIWidDe{0pt}%
\providecommand\NudgeEuBohatstviTopIWidPl{0pt}%
\providecommand\NudgeEuBohatstviTopIWidSk{0pt}%
\begin{figure}[H]
  \centering
  % Editable figure definition for eu_bohatstvi_top1_wid
% Edit freely --- Python will NOT overwrite this file.
% To regenerate defaults, delete this file and re-run the script.
\def\strEuBohatstviTopIWidTitle{Podíl top 1\,\% domácností na čistém jmění (WID, do 2024)}%
\def\strEuBohatstviTopIWidColorbarLabel{podíl top 1\,\% na čistém jmění [\%]}%
\def\strEuBohatstviTopIWidCaption{Podíl top 1\,\% domácností na čistém jmění, mapa \acs{EU}, do~2024. Zdroj dat: \acs{WID}~\cite{wid_world_shweal}.}%
%

% --- Auto-added y-nudge knobs (override with \renewcommand)
\providecommand\NudgeEuBohatstviTopIWidCz{0pt}%
\providecommand\NudgeEuBohatstviTopIWidDk{0pt}%
\providecommand\NudgeEuBohatstviTopIWidAt{6pt}%
\providecommand\NudgeEuBohatstviTopIWidDe{0pt}%
\providecommand\NudgeEuBohatstviTopIWidPl{0pt}%
\providecommand\NudgeEuBohatstviTopIWidSk{0pt}%
\begin{figure}[H]
  \centering
  % Editable figure definition for eu_bohatstvi_top1_wid
% Edit freely --- Python will NOT overwrite this file.
% To regenerate defaults, delete this file and re-run the script.
\def\strEuBohatstviTopIWidTitle{Podíl top 1\,\% domácností na čistém jmění (WID, do 2024)}%
\def\strEuBohatstviTopIWidColorbarLabel{podíl top 1\,\% na čistém jmění [\%]}%
\def\strEuBohatstviTopIWidCaption{Podíl top 1\,\% domácností na čistém jmění, mapa \acs{EU}, do~2024. Zdroj dat: \acs{WID}~\cite{wid_world_shweal}.}%
%

% --- Auto-added y-nudge knobs (override with \renewcommand)
\providecommand\NudgeEuBohatstviTopIWidCz{0pt}%
\providecommand\NudgeEuBohatstviTopIWidDk{0pt}%
\providecommand\NudgeEuBohatstviTopIWidAt{6pt}%
\providecommand\NudgeEuBohatstviTopIWidDe{0pt}%
\providecommand\NudgeEuBohatstviTopIWidPl{0pt}%
\providecommand\NudgeEuBohatstviTopIWidSk{0pt}%
\begin{figure}[H]
  \centering
  \input{../python/figures/eu_bohatstvi_top1_wid.pgf}
  \caption{\strEuBohatstviTopIWidCaption}
  \label{fig:eu_bohatstvi_top1_wid}
\end{figure}

  \caption{\strEuBohatstviTopIWidCaption}
  \label{fig:eu_bohatstvi_top1_wid}
\end{figure}

  \caption{\strEuBohatstviTopIWidCaption}
  \label{fig:eu_bohatstvi_top1_wid}
\end{figure}

